
\begin{abstract}
长江十年禁渔实施后，四大家鱼资源量回升，长江江豚和长江鲟种群亦呈恢复态势，但资源枯竭、江湖连通性受阻、水体污染与外来物种入侵等问题仍较突出。围绕"资源恢复能否稳定转化为生态系统健康改善"核心问题，本文分别描述流域平均食物网、珍稀物种响应和三级食物链的长期行为，再用综合情景评价比较污染与入侵压力下能够保留的禁渔收益。

问题一在流域平均尺度上配对四类基础资源与四大家鱼功能群。以2025年四大家鱼资源量恢复至禁渔前1.8倍作单参数标定，湖北江段单网次渔获量提升120\%留作局地对照。结果显示，2025年四大家鱼总生物量为禁渔前的1.805倍，单位捕捞努力量渔获量($\mathrm{CPUE}^*$)为1.838倍；摄食压力指数升至4.213，基础资源量末值降至0.204。鱼类回升早于底层资源恢复。

问题二把上游资源--鱼群轨迹压缩为综合生态状态指数，用它调节中层猎物承载力，并设置高阻隔、无人工放流和污染情景。基线下长江江豚与长江鲟末值为1.133倍和1.140倍，高阻隔时降至0.901倍和0.856倍；取消人工放流后长江鲟降至0.519倍，污染情景下两者为0.610倍和0.704倍。对照结果显示，通道阻隔同时压低两类物种，放流变化的主要响应落在长江鲟。

长期行为部分采用同一Hastings-Powell三级食物链。参数$K$增大时，代表点$K=1.795$、$2.665$、$4.800$的峰谷差和峰间隔离散程度随之增大，轨道由规则振荡转为复杂非线性波动。短时间窗分析显示，首个连续正李雅普诺夫指数区间位于$K=3.2186$附近；75组扩展测试中，仅5组在$K=3.20\sim3.22$区间保留由负转正的跨越。故本文仅报告"系统动力学复杂性随承载力提高而增强"的方向性结论。

综合评价取积分后期的饵料资源、长江江豚、长江鲟和入侵压力指标，并用熵权与CRITIC平均权重计算情景得分。三种情景依次得到0.662、0.423和0.155；2000次局部参数扰动无一改变排序，稳定率为100\%。这一顺序与"水体污染首先削弱禁渔收益，外来物种入侵进一步导致系统状态恶化"相对应。

十年禁渔后，鱼类资源已经回升，但资源恢复不会自动转化为系统整体健康改善：模型中的底层资源、珍稀物种和长期轨道变化幅度各异。若不同步推进水质治理、江湖连通性修复、珍稀物种定向放流与入侵防控，恢复收益仍可能在向高营养级传递和多重压力叠加时减弱。
\end{abstract}

